%==============================================================================
%  LAMPIRAN
%==============================================================================
\section{Lampiran}

\subsection{Tautan Release Repository GitHub}

Kode sumber lengkap kompiler Arion beserta release Milestone 4 tersedia pada
repositori GitHub berikut.

\begin{itemize}
  \item \textbf{Repositori:} \url{https://github.com/Nusss0/CNA-Tubes-IF2224-2026}
  \item \textbf{Release Milestone 4 (tag \texttt{v0.4.1}):}\\
        \url{https://github.com/Nusss0/CNA-Tubes-IF2224-2026/releases/tag/v0.4.1}
\end{itemize}

\subsection{Pembagian Tugas}

Pembagian tugas dan persentase kontribusi setiap anggota kelompok pada Milestone 4
dirangkum pada Tabel~\ref{tab:tugas}.

\begin{table}[H]
\centering
\caption{Pembagian tugas dan persentase kontribusi Milestone 4.}
\label{tab:tugas}
\renewcommand{\arraystretch}{1.3}
\begin{tabularx}{\textwidth}{@{}l c X c@{}}
\toprule
\textbf{Nama} & \textbf{NIM} & \textbf{Tugas Milestone 4} & \textbf{Kontribusi} \\
\midrule
Stevanus Agustaf Wongso & 13524020 &
Validasi runtime pada stack machine, modul \texttt{RuntimeErrorHook}, deteksi
overflow/underflow, stack corruption \& smashing, runtime type checking
(tantangan bonus). & 25\% \\
\addlinespace
Renuno Yuqa Frinardi & 13524080 &
Inti interpreter / stack machine (\texttt{StackMachine}, \texttt{RuntimeValue},
\texttt{RuntimeContext}), \texttt{TacAdapter}, dan penyusunan kasus uji
Milestone 4. & 25\% \\
\addlinespace
Valentino Daniel Kusumo & 13524104 &
Pembentukan control flow pada generator (\texttt{genIf}, \texttt{genWhile}) 
\& compound block, integrasi pipeline pada \texttt{main}, serta penulisan laporan. & 25\% \\
\addlinespace
Michael James Liman & 13524106 &
Inti \texttt{TacGenerator} (pemetaan alamat, pembentukan ekspresi \& penugasan,
resolusi label, pencetakan), kasus uji ekspresi/control flow, dan penulisan
laporan. & 25\% \\
\midrule
\multicolumn{3}{r}{\textbf{Total}} & \textbf{100\%} \\
\bottomrule
\end{tabularx}
\end{table}

%==============================================================================
%  REFERENSI
%==============================================================================
\section{Referensi}

\begin{enumerate}[label={[\arabic*]}]
  \item A. V. Aho, M. S. Lam, R. Sethi, dan J. D. Ullman, \textit{Compilers:
        Principles, Techniques, and Tools}, edisi ke-2. Boston: Addison-Wesley,
        2006.
  \item Tim Asisten IF2224, \textit{Spesifikasi Milestone 4 -- Tubes IF2224 TBFO:
        Intermediate Code and Interpreter}, Institut Teknologi Bandung, 2026.
  \item Tim Asisten IF2224, \textit{Lampiran Spesifikasi Milestone 4 -- Guidebook
        Konvensi Intermediate Code}, Institut Teknologi Bandung, 2026.
  \item P. Koopman, \textit{Stack Computers: The New Wave}. [Daring]. Tersedia:
        \url{https://users.ece.cmu.edu/~koopman/stack_computers/sec3_2.html}
  \item E. Bendersky, ``Stack frame layout on x86-64,'' 2011. [Daring]. Tersedia:
        \url{https://eli.thegreenplace.net/2011/09/06/stack-frame-layout-on-x86-64}
  \item GEPL, ``Decorated Abstract Syntax Tree,'' University of Minho. [Daring].
        Tersedia: \url{https://epl.di.uminho.pt/~gepl/GEPL_DS/Alma2/dapast.php}
\end{enumerate}
